\section{The odd-cycle law}

Suppose (n) is odd. Every element of (langle\comm\rangle) is an even
power (H^{2k}), whereas (
eceipt=H^n) has an odd exponent. Therefore
\[
\receipt\notin\langle\comm\rangle.
\]
Since (|\comm|=n) is odd and (|\receipt|=2), their subgroups have trivial
intersection. Both lie in the cyclic group (langle H\rangle), so they
commute, and their product has order (2n). Consequently,
\[
\langle H\rangle
=\langle\comm\rangle\times\langle\receipt\rangle
\cong C_n\times C_2
\cong C_{2n}.
\]

The original evolver is recovered from the two receipts. Since
\[
2\left(\frac{n+1}{2}\right)+n=2n+1,
\]
we have
\[
\comm^{(n+1)/2}\receipt
=H^{n+1}H^n
=H^{2n+1}
=H.
\]

\begin{proposition}[Odd receipt splitting]
If (n) is odd, then
\[
\receipt\notin\langle\comm\rangle,
\qquad
\langle H\rangle
=\langle\comm\rangle\times\langle\receipt\rangle,
\qquad
H=\comm^{(n+1)/2}\receipt.
\]
\end{proposition}

At odd cycle length, extension and binary return are independent components
whose product reconstructs evolution.

